\section*{LỜI NÓI ĐẦU}
\thispagestyle{empty}
Em tên là Dương Đình Nhật (20242223M), học viên lớp Kỹ thuật Điện tử (KH) - CH2024B, Trường Điện – Điện tử, Đại học Bách khoa Hà Nội. Em xin trân trọng giới thiệu báo cáo đồ án với đề tài \textbf{“Nghiên cứu, phân tích và tái hiện thực nghiệm thuật toán chuyển giao thông minh ILCHO cho mạng vệ tinh LEO trong 6G NTN”}, dựa trên bài báo gốc M. Choi, M. Park, J. Kim, J.-M. Chung, ``Intelligent Handover Scheme for Improved 6G NTN LEO Satellite Network Performance,'' IEEE Transactions on Mobile Computing, vol. 25, no. 5, 05/2026.

Khác với báo cáo đề xuất trước đó khảo sát ba hướng tiếp cận AI cho handover ở mạng mặt đất 5G, trong báo cáo này em đi sâu vào một bài toán cụ thể hơn và mới hơn: chuyển giao (handover) cho người dùng trong mạng phi mặt đất (Non-Terrestrial Network - NTN) dùng chòm vệ tinh quỹ đạo thấp (LEO) mật độ cao — bối cảnh then chốt của 6G. Em đã đọc, phân tích chi tiết bài báo ILCHO, đồng thời tự xây dựng lại toàn bộ mã nguồn (bài báo không công khai mã nguồn) để tái hiện và đối chiếu số liệu công bố, từ đó đánh giá một cách trung thực phần nào tái hiện được và phần nào chưa.

Để có thể hoàn thành đồ án này, chắc chắn không phải chỉ là kết quả được thực hiện trong một kỳ làm đồ án, mà đó là cả quá trình tích lũy kiến thức của em tại trường Điện - Điện tử, Đại học Bách khoa Hà Nội. Em xin gửi lời cảm ơn chân thành đến các thầy cô đã tận tình giảng dạy, hướng dẫn em trong suốt quá trình học tập vừa qua. Các thầy cô không chỉ cho em kiến thức bổ ích mà còn đặt cho em một nền tảng tư duy vững chắc, làm hành trang quý báu trong công việc và cuộc sống của em sau này.

Em xin gửi lời cảm ơn đặc biệt tới \textbf{TS. Phùng Thị Kiều Hà} đã cho em cơ hội làm đề tài và trực tiếp hướng dẫn, chỉ bảo em trong suốt thời gian làm đồ án này. Xin chân thành cảm ơn những lời khuyên và định hướng giúp đỡ của cô đã giúp em mở mang không chỉ là kiến thức mà còn thái độ, tinh thần làm việc nghiêm túc, trung thực với số liệu để hoàn thành đồ án một cách hiệu quả.

\cleardoublepage
